% =====================================================================
% Section 2: Related Work
% =====================================================================
%
% Status: Draft §2 Related Work v7 (2026-05-19)
%   v7: half-page advisor-driven rewrite.
%       Three buckets (Standard RAG / Graph RAG / Evidence-enhanced
%       RAG), each ending with an explicit limitations sentence.
%       Closing Summary keeps the §1 conversion-problem tie-back.
%   v6: see archive/02_related_v6_pre-half-page-rewrite_20260519.tex
%
% =====================================================================

\section{Related Work}
\label{sec:related}

\paragraph{Standard RAG.}
Standard retrieval-augmented generation methods rank passages by
lexical or dense similarity to the question, with BM25
\citep{robertson2009bm25} and DPR \citep{karpukhin2020dpr} as the
canonical sparse and dense baselines. RAPTOR
\citep{sarthi2024raptor} addresses limited cross-passage context by
recursively clustering and summarizing passages into a retrieval
tree. Iterative chain-of-thought retrievers such as IRCoT
\citep{trivedi2023ircot} and Self-Ask \citep{press2023selfask}
drive multi-hop retrieval through language-model-generated reasoning
steps, alternating between generation and lookup. \emph{Limitation.}
These systems score each passage independently or follow flat
neighbor relations, so the retrieval state remains a list or beam
of passages without an explicit cross-document structure; bridge
passages whose surface similarity to the question is low are
systematically underweighted.

\paragraph{Graph RAG.}
A second line of work organizes the corpus into an explicit graph
to expose cross-document connections. HippoRAG~2
\citep{gutierrez2025hipporag} extracts OpenIE facts and runs
personalized PageRank over a corpus-wide passage--phrase graph
seeded from the question; PropRAG \citep{wang2025proprag} replaces
triples with propositions and discovers multi-hop evidence by beam
search; GraphRAG \citep{edge2024graphrag} builds community-summary
graphs over the corpus; HGRAG \citep{hgrag2024} models entities and
passages within a hypergraph and propagates relevance across
granularities. Recent variants explore hierarchical multi-granular
graphs \citep{hu2026mgrangrag}, sentence-level graphs
\citep{liang2026sentgraph}, and CPU-only PageRank-style retrieval
\citep{wang2026sprig}, while earlier graph-based readers integrate
evidence via graph neural message passing
\citep{qiu2019dfgn,song2022evidenceintegration}.
\emph{Limitation.} These systems condition on a pre-built global
structure, so question conditioning enters only through seed
distributions or post-hoc search; the resulting node-level
relevance does not directly translate into a short reader-facing
passage list that is both connected across documents and covering
of the entities and relations the question asks about.

\paragraph{Evidence-enhanced RAG.}
A third line uses graph or KG signals to refine the evidence that
reaches the reader. NeocorRAG \citep{peng2026neocorrag} mines
\emph{evidence chains} from retrieved text via activation-based
candidate search and constrained decoding, using the chains to
filter and re-present passages. Relink \citep{huang2026relink}
addresses KG incompleteness by reconstructing query-specific
evidence graphs and repairing missing facts at query time. ToG
\citep{sun2024tog} performs LLM-guided traversal over a pre-built
knowledge graph, while KAPING \citep{baek2023kaping} retrieves
subgraphs from a KG and prompts a language model with the
subgraph as context. \emph{Limitation.} These methods operate
\emph{on top of} an underlying retriever or KG: chain mining
inherits any miss made by the upstream retriever, KG-based
reconstruction depends on KG completeness, and per-question
LLM-guided traversal is expensive at scale.

\paragraph{Summary.}
Standard RAG, Graph RAG, and Evidence-enhanced RAG sharpen
respectively passage-level similarity, corpus-level graph
structure, and post-retrieval evidence selection. None of the three
takes the conversion of graph relevance into a connected and
covering reader-facing passage list as its central object.
\methodname{} addresses this conversion within the GraphRAG family
by combining query-local propagation over a fact-source substrate
with evidence-need mining and coverage-aware reranking on top of
an offline OpenIE graph.
